\section{Balanced Weil classes}\label{sec:weil}

Let $A/\Qbar$ have good reduction at the fixed place $w$, and let a CM
field $E$ act on $A$ up to isogeny.  We write $e\mapsto\bar e$ for its
complex conjugation.  Suppose that a polarization $\lambda$ induces this
involution on $E$.  Put
\[
 V=H_1(A(\C),\Q),\qquad \dim_E V=2s,\qquad
 g=\dim A=s[E:\Q],
\]
where $s\geq 1$.  We assume that the Hodge multiplicities at every
embedding of $E$ are $(s,s)$.  The polarization is an alternating form
$\psi:V\times V\longrightarrow\Q(1)$ satisfying
\[
 \psi(ex,y)=\psi(x,\bar e y).
\]

\begin{definition}\label{def:weil-space}
The rational space of \emph{Weil classes} associated with $E$ is the
subspace
\[
 W_E^s(A)\subset
 \bigwedge\nolimits_{\Q}^{2s} V^\vee\otimes\Q(s)
       =H^{2s}(A(\C),\Q(s))
\]
whose scalar extension to a splitting field $\Omega$ of $E$ is
\begin{equation}\label{eq:weil-pure-space}
 W_E^s(A)\otimes_\Q\Omega
 =\bigoplus_{\tau:E\hookrightarrow\Omega}
       \bigwedge\nolimits_\Omega^{2s}V_\tau^\vee
          \otimes\Q(s)_\Omega.
\end{equation}
Here each summand is extended by zero on the other eigenspaces.  The
right side is Galois stable, so it defines the indicated rational
subspace.
\end{definition}

The balanced multiplicities imply that every element of $W_E^s(A)$ is
a Hodge class: each determinant in \eqref{eq:weil-pure-space} has
untwisted type $(s,s)$.  We use its compatible realizations as an
absolute Hodge class; see \cite{Deligne1982}.  The same eigenspace
description of its support holds after specialization, since the
$E$-action specializes compatibly.

Our goal is the Lefschetz pairing property for these specialized
classes.  In fact, the reduction satisfies a dichotomy determined by
its rational endomorphisms.  Let $\dagger$ be Rosati for the reduced
polarization and put
\begin{equation}\label{eq:weil-U}
 \mathcal U=
 \{u\in\End^0(A_0):u^\dagger=u,\quad
                      ue=\bar e u\ \text{for every }e\in E\}.
\end{equation}

\begin{proposition}\label{prop:weil}
Every $\gamma\in W_E^s(A)$ satisfies the Lefschetz pairing property
at $w$.  More precisely, if $\mathcal U$ in \eqref{eq:weil-U} has no
invertible element, all complementary pairings are zero.  If it has
an invertible element, $\gamma_0$ is a simultaneous rational
Lefschetz class.
\end{proposition}

We first prove the vanishing alternative.  For the other alternative,
an invertible $u$ supplies a divisor whose form is nondegenerate on
each $E$-eigenspace separately.  We will lift this divisor and use it
to obtain a rational identity for the determinant classes.

\subsection{The orthogonality alternative}

\begin{lemma}\label{lem:weil-orthogonality}
If $\mathcal U$ contains no invertible element, every specialization of
a class in $W_E^s(A)$ is orthogonal to all complementary Lefschetz
classes, in each cohomological realization.
\end{lemma}

\begin{proof}
We prove the contrapositive.  A nonzero pairing will produce a divisor
form nondegenerate on one $E$-eigenspace.  We will then show that a
rational endomorphism in $\mathcal U$ with this property is invertible.

Fix one realization and extend its coefficient field to an
algebraically closed field $\Omega$.  At the crystalline component we
also include the coefficient extension used for the specialized
class.  In this proof, $V_\tau$ also denotes the corresponding
eigenspace in the dual of first cohomology of $A_0$, identified by
specialization and comparison.  Write $L^j_\Omega(A_0)$ for the $\Omega$-span of products of
$j$ divisor classes.  Since $[E:\Q]\geq2$, we have $g-2s\geq0$.
Hard Lefschetz on the Lefschetz algebra gives
\begin{equation}\label{eq:weil-hard-lefschetz}
 L^{g-s}_\Omega(A_0)
        =\psi_0^{g-2s}L^s_\Omega(A_0).
\end{equation}
The restriction of hard Lefschetz to this algebra, including for
crystalline cohomology, follows from the Lefschetz correspondences
of \cite[Theorem~5.9]{Milne1999}.

Suppose a Weil class pairs nontrivially with
$D_1\cdots D_s\psi_0^{g-2s}$, and choose a pure term, supported on
$V_\tau$, whose contribution is nonzero.  This term already has
degree $2s$ on $V_\tau$, so no further factor in a nonzero top-degree
product can use that space.  A factor from $\psi_0$ which uses
$V_{\bar\tau}$ also uses $V_\tau$; hence the polarization factors
contribute to neither space in this pair.  All $2s$ degrees on
$V_{\bar\tau}$ must therefore come from the $s$ divisor factors.
Each has degree two, so each must contribute its restriction to
$V_{\bar\tau}$, and those restrictions have nonzero product.
The coefficient of $t_1\cdots t_s$ in
\[
 (t_1D_1+\cdots+t_sD_s)^s\big|_{V_{\bar\tau}}
\]
is $s!$ times that product.  Some $\Omega$-linear combination of the
divisors therefore restricts to a nondegenerate alternating form on
$V_{\bar\tau}$.

We identify which endomorphisms can give these restrictions.  Let
$p_\xi\in E\otimes_\Q\Omega$ be the eigenspace projectors.  Rosati
satisfies $p_\xi^\dagger=p_{\bar\xi}$.  A divisor corresponds under
$\psi_0$ to a Rosati-symmetric endomorphism $t$.  Define
\begin{equation}\label{eq:weil-block-projection}
 \Pi(t)=\sum_\xi p_{\bar\xi}t p_\xi.
\end{equation}
This operation takes place inside
$\End^0(A_0)\otimes_\Q\Omega$.  Moreover,
\[
 (p_{\bar\xi}t p_\xi)^\dagger
       =p_{\bar\xi}t^\dagger p_\xi,
 \qquad
 \Pi(t)e=\bar e\Pi(t).
\]
Thus $\Pi(t)\in\mathcal U\otimes_\Q\Omega$.  Conversely, these are
exactly the blocks allowed by the second equation in
\eqref{eq:weil-U}.  For $x,y\in V_\xi$, the form
$\psi_0(x,ty)$ depends only on $p_{\bar\xi}t p_\xi$.
It follows that projecting by \eqref{eq:weil-block-projection}
preserves all the same-embedding restrictions in question.

Nondegeneracy on $V_{\bar\tau}$ is a polynomial open condition on
$\mathcal U\otimes_\Q\Omega$, and we have shown this open is
nonempty.  The rational points of the rational vector space
$\mathcal U$ are Zariski dense after extension to $\Omega$.  Hence
there is a rational $u\in\mathcal U$ for which
$\psi_0(\cdot,u\cdot)$ is nondegenerate on $V_{\bar\tau}$.

To prove invertibility, we use the fact that a nonzero $E$-stable
abelian subvariety has a nonzero homology space at every embedding
of $E$.  Suppose that $u$ is not invertible, and choose an integer
$N>0$ such that $Nu$ is an endomorphism.  The reduced identity
component $C=((\ker(Nu))^0)_{\mathrm{red}}$ is a
positive-dimensional abelian subvariety.  The relation $ue=\bar e u$ makes $C$ stable under
the rational action of $E$.  This action is unital and therefore
faithful.  Choose a primitive element $e$ of $E$ and an integer
$M>0$ such that $Me$ acts by an integral endomorphism of $C$.
The exterior algebra description of cohomology and normalized
top-degree pullback give, in every realization
\cite[Appendix, Proposition~(A.2)(a),(b)]{Milne1999},
\[
 \det\bigl(n-e\mid H^1(C)\bigr)
    =M^{-2\dim C}\deg\bigl([Mn]-Me\bigr)
\]
for every integer $n$ for which the map on the right is an
isogeny.  All but finitely many integers have this property.
The right side is rational and independent of the realization.
Polynomial interpolation therefore gives one characteristic
polynomial in $\Q[T]$, identical also in crystalline cohomology.
Its roots lie among
the conjugates of $e$, and rationality forces their multiplicities to
be equal.  Since $C\ne0$, every embedding of $E$ occurs with positive
multiplicity.  The same statement holds on first homology, in the
chosen realization. Rational Poincar\'e reducibility
\cite[Proposition~12.1]{MilneAV} makes the inclusion
\(C\hookrightarrow A_0\) injective on first homology in every
realization. This inclusion is \(E\)-equivariant, and its
image is killed by \(u\). Thus the nonzero \(\bar\tau\)-eigenspace
of \(H_1(C)\) shows that the kernel of \(u\) meets
\(V_{\bar\tau}\) nontrivially.  Such a vector lies in the radical of
$\psi_0(\cdot,u\cdot)|_{V_{\bar\tau}}$, a contradiction.

A nonzero pairing would therefore give an invertible element of
$\mathcal U$.  Taking the contrapositive proves the lemma.  The
argument applies separately to every realization, so all the
pairings in this case have the common rational value zero.
\end{proof}

\subsection{A divisor identity on a rational summand}

For the rest of the section, suppose that $u\in\mathcal U$ is
invertible.  It gives a rational divisor $D^u$ on $A_0$ with
alternating form
\begin{equation}\label{eq:weil-extra-divisor}
 D^u(x,y)=\psi_0(x,uy).
\end{equation}
This form is nondegenerate and pairs each embedding of $E$ only with
itself: $u$ exchanges opposite eigenspaces, whereas $\psi_0$ pairs
opposite eigenspaces.

Prescribed-filtration transport will replace $A_0$ by a power and
then by an isogenous special fiber $B_0$.  We will retain the first
factor through a rational projector on $B_0$, and lift that projector,
the $E$-action on its image, and the transported divisor to $B$.
The next lemma explains the resulting divisor identity: a divisor
nondegenerate separately on those eigenspaces generates their
determinant classes.  A rational projector here means an idempotent
in the rational endomorphism algebra.

\begin{lemma}\label{lem:weil-divisor-span}
Let $B$ be a complex abelian variety, let $R\in\End^0(B)$ be a
rational projector, and suppose that there is an algebra action
$E\longrightarrow R\End^0(B)R$ whose unit maps to $R$.
Let $T=\operatorname{im}(R)$ in rational first homology, with
$\dim_E T=2s$; the indicated endomorphisms act by zero on
$\ker(R)$.  Suppose there is a rational divisor class $D$ on
$B$ such that $R^*D=D$ and, after splitting $E$, its alternating form
is
\[
 D=\sum_\tau D_\tau,
 \qquad D_\tau\in\bigwedge\nolimits^2 T_\tau^\vee\otimes\Q(1),
\]
with each $D_\tau$ nondegenerate.  Then the rational subspace whose
scalar extension is
\[
 \bigoplus_\tau\bigwedge\nolimits^{2s}T_\tau^\vee\otimes\Q(s),
\]
extended by zero on the complementary summand, consists of rational
Lefschetz classes on $B$.
\end{lemma}

\begin{proof}
For $b\in E$, let $\widetilde b\in\End^0(B)$ act by $b$ on $T$ and
by zero on the complement.  Its pullback of a rational divisor is a
rational divisor: clear a denominator in $\widetilde b$, and use
that multiplication by $N$ acts as $N^2$ on $H^2$.  Over a splitting
field $\Omega$, choosing Tate-line bases, we have
\[
 (\widetilde b^*D)^s
       =\left(\sum_\tau\tau(b)^2D_\tau\right)^s.
\]
The coordinates $x_\tau=\tau(b)$ are independent coordinates on the
scalar extension of the rational vector group underlying $E$.
In the corresponding vector-valued polynomial
\begin{equation}\label{eq:weil-coefficient-polynomial}
 P((x_\tau)_\tau)=\left(\sum_\tau x_\tau^2D_\tau\right)^s,
\end{equation}
the coefficient of $x_\tau^{2s}$ is $D_\tau^s$.  Nondegeneracy makes
this coefficient nonzero, so it spans
$\bigwedge^{2s}T_\tau^\vee\otimes\Q(s)_\Omega$.

Let $M$ be the rational span of the Lefschetz classes
$(\widetilde b^*D)^s$ for $b\in E$.  The set $E$ of rational points
is Zariski dense in its vector group, also after extension to
$\Omega$.  Any linear functional annihilating $M\otimes_\Q\Omega$
annihilates all evaluations of \eqref{eq:weil-coefficient-polynomial}
on those points.  It therefore annihilates the polynomial, and hence
each coefficient.  Every $D_\tau^s$ belongs to
$M\otimes_\Q\Omega$.

This proves the required inclusion after scalar extension.  Both
subspaces are rational; applying faithful scalar extension to their
quotient proves the inclusion over $\Q$.  In particular it gives
rational coefficients for every rational class in the asserted
subspace, without requiring an individual embedding line to be
defined over $\Q$.
\end{proof}

\subsection{A lift preserving the divisor}

We now construct the auxiliary variety to which
Lemma~\ref{lem:weil-divisor-span} applies.  To transport the Weil class,
we use the $E$-linear polarization similitudes whose determinant has
the same transformation law as a codimension-$s$ Tate twist.
Choose a rational basis
for the Betti Tate line, let $c$ be the multiplier of
$\GSp(V,\psi)$, and set
\begin{equation}\label{eq:weil-group}
 G=\{a\in\GSp(V,\psi)\cap\GL_E(V):
                         \det_E(a)=c(a)^s\}.
\end{equation}
On points over a $\Q$-algebra $R$, the determinant equality is in
$(E\otimes_\Q R)^\times$, with the right side embedded diagonally.

\begin{lemma}\label{lem:weil-group}
The Hodge homomorphism of $A$ factors through $G$, and its conjugacy
class defines a Shimura datum of Hodge type through the displayed
embedding.  Over an algebraically closed coefficient field, the orbit
of the corresponding filtration on $H^1$ consists exactly of the
$E$-stable Lagrangians for the dual polarization pairing whose
dimension at each embedding of $E$ is $s$.  Using a basis of its Tate
line, every alternating form in \eqref{eq:weil-pure-space} transforms
by $c^s$.
\end{lemma}

\begin{proof}
The polarization pairs $V_\tau$ with $V_{\bar\tau}$, and no other pair
of eigenspaces.  Over an algebraically closed field, choose one
embedding in each conjugate pair.  An element of $G$ is specified by
$c$ and matrices $a_\tau\in\GL_{2s}$ with
$\det(a_\tau)=c^s$; the matrices on the opposite spaces are then
determined by the polarization.  In particular,
\[
 \ker(c:G\longrightarrow\mathbb G_m)
       \simeq\prod_{\{\tau,\bar\tau\}}\SL_{2s}.
\]
The multiplier is surjective geometrically.  Thus $G$ is connected
reductive, with the same adjoint group as the group of $E$-linear
polarization similitudes.

The two Hodge characters occur $s$ times at every embedding.  Their
product is the polarization multiplier, so the Hodge homomorphism
satisfies the determinant condition in \eqref{eq:weil-group}.  Its
weight is rational and central.  On the adjoint Lie algebra the Hodge
types are among $(-1,1),(0,0),(1,-1)$, as follows from the weight-one
representation.  Conjugation by the image of $i$ is a Cartan
involution, by the polarization.  At every real embedding of the
maximal totally real subfield of $E$, the adjoint group has unitary
signature $(s,s)$.  Its Hodge homomorphism is nontrivial there, so no
rational simple adjoint factor has trivial Hodge homomorphism.
These are the Shimura-datum conditions for the given Hodge embedding.

On $H^1$ the polarization still pairs opposite eigenspaces.  An
$E$-stable Lagrangian is specified by an $s$-dimensional subspace at
one embedding of each pair; the subspace at the opposite embedding
is its annihilator.  The factor $\SL_{2s}$ acts transitively on these
subspaces, as in the cohomological orbit calculation of
Theorem~\ref{thm:filtration}.  Finally, an alternating determinant form
evaluated on $a_\tau v_1,\ldots,a_\tau v_{2s}$ is multiplied by
$\det(a_\tau)=c^s$.  This proves the asserted transformation law.
\end{proof}

\begin{proof}[Proof of Proposition~\ref{prop:weil}]
The vanishing alternative is Lemma~\ref{lem:weil-orthogonality}.
Suppose that $u$ is invertible, and use the divisor $D^u$ from
\eqref{eq:weil-extra-divisor}.

\medskip\noindent
\emph{A filtration preserving the divisor.}

Let $\mathcal D=H^1_{\mathrm{cris}}(A_0/W(\F))[1/p]$, with its
contravariant convention, and let $\beta$ be the dual polarization
pairing, taking values in the isocrystal line of slope one.  The
actions on $\mathcal D$ satisfy
\[
 \beta(ux,y)=\beta(x,uy),\qquad
 \beta(ex,y)=\beta(x,\bar e y),\qquad ue=\bar e u.
\]
All these endomorphisms and the polarization have finite field
descent.  Theorem~\ref{thm:filtration} supplies, over a finite
extension with finite local field descent, a weakly admissible
filtration $\Fil^1\mathcal D=L$ which is $E$- and $u$-stable,
is Lagrangian for $\beta$, and has dimension $s$ at every embedding
of $E$.  By Lemma~\ref{lem:weil-group} it lies geometrically in the
required $G$-orbit.  The
\hyperref[par:transport-tensor-descent]{finite-descent argument} in
Proposition~\ref{prop:transport} shows that its full finite family of
Hodge-type structure tensors descends from the original number-field
point by crystalline comparison.  Taking a compositum with the local
field of definition of $L$ gives the common finite local-field descent
required by its prescribed-filtration clause.

\medskip\noindent
\emph{Transport and the first rational summand.}
Apply Proposition~\ref{prop:transport} to $G$ and $\gamma$ with this
filtration.  It gives a lift $B$, a
quasi-isogeny $a:A_0^n\longrightarrow B_0$, and a rational Hodge
class $\kappa$ on $B$, with
\[
 f_0=aj:A_0\longrightarrow B_0,\qquad
 f_0^*\kappa_0=\gamma_0,
\]
simultaneously in the realizations under consideration.  The
filtration on $B$, marked by $a$, is the direct sum of $n$ copies of
the chosen filtration.  Here $j$ includes the first factor.  The map
$f_0$ identifies $A_0$ with a rational direct summand of $B_0$.
Its retraction $q_0$ and the corresponding projector $R_0$ are
\begin{equation}\label{eq:weil-projector-identity}
 q_0=\operatorname{pr}_1a^{-1},\qquad R_0=f_0q_0.
\end{equation}
Thus $q_0f_0=\id$ and $R_0^2=R_0$.  Pulling back
$f_0^*\kappa_0=\gamma_0$ by $q_0$ gives
\[
 R_0^*\kappa_0=q_0^*\gamma_0.
\]
Act by $E$ on $\operatorname{im}(R_0)$ using $e\mapsto f_0e q_0$,
and by zero on the complement.  Also put $D_0=q_0^*D^u$.

\medskip\noindent
\emph{Lifting the projector, action, and divisor.}
The projector and this $E$-action preserve the repeated filtration.
By Corollary~\ref{cor:filtration-divisor}, $D^u$ satisfies the divisor
lifting condition for the chosen filtration on the first summand.
The crystalline map induced by $q_0=\operatorname{pr}_1a^{-1}$
respects that filtration and the marked repeated filtration on $B_0$.
Hence $D_0=q_0^*D^u$ satisfies the same condition on $B$.
Apply Lemma~\ref{lem:lift-hom} to \(R_0\), to a finite
\(\Q\)-basis of the action of \(E\) on \(\operatorname{im}(R_0)\), and to the
symmetric homomorphism \(h_{D_0}:B_0\to B_0^\vee\) defined by \(D_0\).
Denote their lifts by \(R\), the corresponding basis of maps, and
\(h_D\). Injectivity of specialization preserves \(R^2=R\), the
rational multiplication table of that basis, and its unit \(R\).
It therefore gives an algebra action \(E\to R\End^0(B)R\).
The symmetric map \(h_D\) defines a rational divisor class \(D\).
Moreover, \(q_0R_0=q_0\), so
\[
 R_0^\vee h_{D_0}R_0=h_{D_0}.
\]
Injectivity lifts this identity to \(R^\vee h_DR=h_D\), which is
exactly \(R^*D=D\). Thus the required projector, action, and divisor
all exist on the one auxiliary variety $B$.

\medskip\noindent
\emph{The rational identity and its specialization.}
Consider now rational Betti cohomology at a complex realization of
$B$.  The summand $T=\operatorname{im}(R)$ has dimension $2s$ at
each embedding of $E$.  The form of $D$ is supported on this summand
and pairs each embedding with itself nondegenerately.  Also
$R^*\kappa$ is supported on the pure determinant spaces of this
summand.  These assertions can be checked after extension to any
one prime-to-$p$ realization: for the divisor they follow from its
specified specialization, and for the Hodge class they follow from
the equality $R_0^*\kappa_0=q_0^*\gamma_0$ and the
pure support of $\gamma_0$.  Compatibility of the transported
realizations then gives these support and nondegeneracy assertions
in Betti cohomology.

Lemma~\ref{lem:weil-divisor-span} now expresses the rational Betti
class $R^*\kappa$ as a finite rational linear combination of products
of rational divisor classes on $B$.  The realization and
specialization compatibilities in
Proposition~\ref{prop:transport}, together with those for divisors
and homomorphisms, specialize this identity with the same rational
coefficients in every component.  Hence $R_0^*\kappa_0$ is a
simultaneous rational Lefschetz class.  Finally,
\[
 \gamma_0=f_0^*R_0^*\kappa_0,
\]
because $R_0f_0=f_0$.  Pullback by a rational homomorphism preserves
simultaneous rational Lefschetz classes, so the same is true of
$\gamma_0$.  Its complementary traces are therefore one common
rational intersection number, proving the second alternative and
the proposition.
\end{proof}
